\section*{Part E. Independent certificates and the CM Selmer deductions}

Put \(K=\mathbb Q(i)\), \(p=5\), and
\(W_{m,d}=\operatorname{Ind}_{G_K}^{G_{\mathbb Q}}\psi_d^{-m}\),
with the homological convention
\(\psi_d\overline\psi_d=\chi_5\). Thus \(W_{m,d}\) is the
cohomological representation of the weight \(m+1\) CM form,
without an additional Tate twist. The calculation below
separates the numerical certificates from the Iwasawa input.

\begin{statement}{E.3}{Proved}{The six balanced lines by a class-group calculation}
For \(t=1,2,3\),
\[
 H_f^1(\mathbb Q,\mathbb Q_5(-t))
 =H_f^1(\mathbb Q,\mathbb Q_5(-t)\epsilon_K)=0.
\]
This proof uses neither modular symbols nor Kato's theorem,
nor a conjectural nonvanishing statement.
\end{statement}
\begin{proof}
Let \(L=\mathbb Q(\zeta_{20})=K(\zeta_5)\).
We first give a small class-number certificate. The polynomial
\(f(X)=X^8-X^6+X^4-X^2+1\) has discriminant \(4000000\),
and
\[
 \operatorname{Res}_X(f(X),X^3-X+1)=5.                 \tag{E.3.1}
\]
Both equalities are exact integer computations in
\texttt{code/de\_balanced.py}. Minkowski's bound is
\[
 \frac{8!}{8^8}\left(\frac4\pi\right)^4\sqrt{4000000}
 < 12.6321 <13.
\]
For the rational primes \(2,3,7,11\), the least residue degrees
in \(L\) are respectively \(4,4,4,2\); their prime-ideal norms
are all greater than twelve. These assertions are also verified
by factoring \(f\) modulo each of these primes. At five there
are two primes, both of norm five, since five has ramification
index four and splits in \(K\). Equation (E.3.1) makes one
of them principal, and its Galois conjugate is principal as
well. These are the only prime ideals of norm below thirteen.
Every ideal class therefore has a principal representative by
Minkowski, so the class number of \(L\) is one.

Let \(S\) consist of the places above \(2,5\), together with
the infinite places, and put \(A=\operatorname{Gal}(L/K)\).
Then \(A\) is cyclic of order four. There is one prime above
two and there are two above five; each is fixed by \(A\).
The Brauer invariant-sum sequence and the class-number
calculation give, equivariantly,
\[
 H^2(G_{L,S},\mu_5)
 =\operatorname{Br}(\mathcal O_{L,S})[5]
 =\ker(\mathbb F_5^3\xrightarrow{\sum}\mathbb F_5)
 \simeq\mathbb F_5^2,                                  \tag{E.3.2}
\]
with trivial \(A\)-action. The first equality follows from
the Kummer exact sequence because the ideal class group
vanishes. The action on local Brauer invariants is trivial;
no choice of roots of unity alters it.

Let \(\omega:A\to\mathbb F_5^\times\) be the cyclotomic
character. It is an isomorphism. The action on
\(H^2(G_{L,S},\mathbb F_5(j))\) is the action in (E.3.2)
tensored by \(\omega^{j-1}\). For \(j=2,3,4\) this character
is nontrivial. Since four is invertible modulo five,
restriction and averaging give
\[
 H^2(G_{K,S},\mathbb F_5(j))=0\qquad(j=2,3,4).           \tag{E.3.3}
\]
The extensions
\(0\to\mathbb F_5(j)\to\mathbb Z/5^n(j)
\to\mathbb Z/5^{n-1}(j)\to0\) prove the same vanishing
at every finite exponent by induction. Finite-ramification
cohomology with finite coefficients is finite, so its
inverse systems are Mittag--Leffler. Jannsen's comparison
then gives
\[
 H^2(G_{K,S},\mathbb Q_5(j))=0\qquad(j=2,3,4).           \tag{E.3.4}
\]
The finiteness and comparison inputs are
\cite[I, Corollary 4.15]{DE_MilneADT} and
\cite[Proposition 1.6, Theorem 2.2]{DE_Jannsen1988}.

Take \(V=\mathbb Q_5(-t)\) over \(K\). Its local finite
condition is zero everywhere. At a place above five,
\(D_{\rm dR}(V)/\operatorname{Fil}^0=0\) and
\(H^0(K_v,V)=0\); the crystalline Bloch--Kato dimension
formula \cite[Corollary 3.8.4]{DE_BlochKato1990} therefore gives \(H_f^1(K_v,V)=0\). Away from five,
the unramified local group is the cokernel of
\(\operatorname{Frob}_v-1\), which is invertible on \(V\)
because its eigenvalue is \((\#k_v)^{-t}\ne1\).
Complex places contribute zero. Poitou--Tate duality makes
the kernel of global localization in \(H^1(G_{K,S},V)\)
a quotient of \(H^2(G_{K,S},V^*(1))^*\). That last group
is zero by (E.3.4), since \(V^*(1)=\mathbb Q_5(t+1)\).
Thus \(H_f^1(K,V)=0\).

Finally restriction from \(\mathbb Q\) to \(K\) is injective
on rational cohomology for each of the two characters
\(\mathbb Q_5(-t)\) and \(\mathbb Q_5(-t)\epsilon_K\):
corestriction composed with restriction is multiplication
by two. Restriction preserves the finite local conditions,
and both restricted characters equal \(V\). This proves
all six asserted vanishings.
\end{proof}

\begin{statement}{E.2}{Proved}{Independent verification of all twenty-four CM certificates}
For \(d=1,6\), \(1\le m\le6\), \(t\ge0\), and
\(m+2t\le6\), the twenty-four modular-symbol certificates
of LM8 F.16--F.18 are correct, and each decides nonvanishing
at the stated precision. Kato's Theorem 17.4(1)--(2) applies
to the associated ordinary CM form. With its cohomological
normalization, the resulting conclusion is
\[
 H_f^1(\mathbb Q,W_{m,d}(-t))=0.                        \tag{E.2.1}
\]
The proof includes the specialization, restriction-kernel,
and integral-lattice steps. It does not claim the independent
Rubin computation requested in Task E.1.
\end{statement}
\begin{proof}
We specify both the computation and its arithmetic interpretation.
Let \(g_m\) be the theta series of the \(m\)th power of the
Hecke character of \(y^2=x^3-x\). An ideal prime to two has
a unique Gaussian generator congruent to one modulo
\((1+i)^3\); this describes that character explicitly.
Its \(m\)th power has conductor \((1+i)^3\) for odd \(m\),
\((1+i)^2\) for \(m\equiv2\pmod4\), and conductor one
for \(m\equiv0\pmod4\). The associated primitive theta
form has weight \(m+1\), level four times the conductor
norm, and character \(\epsilon_K^{m+1}\). Its
cohomological representation restricts to
\(\psi_1^{-m}\oplus\overline\psi_1^{-m}\), so it is
\(W_{m,1}\). For odd \(m\), the form for \(d=6\) is
its twist by the primitive quadratic character \(\chi_{24}\);
for even \(m\) it is unchanged. In particular five divides
none of these levels.

The prime-five traces are
\((-1+2i)^m+(-1-2i)^m\). Thus the complete local input is
\[
\begin{array}{c|rrrrrr}
m&1&2&3&4&5&6\\\hline
N_m&32&16&32&4&32&16\\
a_5&-2&-6&22&-14&-82&234\\
\alpha_m\bmod25&13&19&22&11&18&9
\end{array}                                                   \tag{E.2.2}
\]
Here \(\alpha_m\) is the unique unit root of
\(X^2-a_5X+5^m\). Since \(\chi_{24}(5)=1\), the roots
are the same for the two curves. These checks prove good
ordinarity in every case, including weights six and seven;
no bound comparing the weight with five is required.

Put \(w=m-1\). We use polynomial Manin generators with
bottom rows
\((1,d)\), \(0\le d<N_m\), and \((c,1)\), \(c\) even
modulo \(N_m\), and the monomials \(X^{w-j}Y^j\).
For a determinant-one matrix \(h\) the right Manin action
is \([P,g]h=[P\circ h,gh]\). Scaling the bottom row by
\(u\) multiplies it by \(\epsilon_K(u)^{m+1}\).
The relations are \(1+S\) and \(1+R+R^2\), for
\[
 S=\begin{pmatrix}0&-1\\1&0\end{pmatrix},\qquad
 R=\begin{pmatrix}0&-1\\1&-1\end{pmatrix}.
\]
The Hecke cosets are
\(\left(\begin{smallmatrix}1&a\\0&\ell\end{smallmatrix}\right)\)
and \(\left(\begin{smallmatrix}\ell&0\\0&1\end{smallmatrix}\right)\),
using the integral adjugate on polynomial coefficients and
the nebentypus on the last coset. These are the general
Manin-symbol conventions of \cite{DE_Stein2007}.

The new program \texttt{code/de\_cm\_audit.py} imports
no supplied program. It takes only the explicit integer
columns \(\Phi\) from the six attached certificates, recorded
in \texttt{de\_cm\_input.json} with their source hashes.
It verifies all Manin relations and both Hecke equations
\(T_3\Phi=0\), \(T_5\Phi=a_5\Phi\) exactly over
\(\mathbb Z\). It independently constructs their complete
coefficient matrix and computes its rank over \(\mathbb F_5\).
The outputs are
\[
\begin{array}{c|rrrrrr}
m&1&2&3&4&5&6\\\hline
D&48&48&144&24&240&144\\
\operatorname{rank}_{\mathbb F_5}(\text{Manin relations})
 &39&40&120&20&200&120\\
\operatorname{rank}_{\mathbb F_5}(\text{all equations})
 &46&46&142&22&238&142\\
\operatorname{rank}_{\mathbb F_5}\Phi&2&2&2&2&2&2
\end{array}                                                   \tag{E.2.3}
\]
These checks prove the full integral lattice assertion.
Indeed an arbitrary integral solution reduces modulo five
to a combination of the two columns. Subtract that
combination, divide by five, and repeat. Completeness gives
a \(\mathbb Z_5\)-combination of \(\Phi\). Conversely the
columns satisfy every equation. Thus there is no unverified
five-power index. The CM form supplies its two independent
period functionals in this kernel; (E.2.3) shows that they
exhaust it. Only rational period identifications are needed.

For clarity the independent path algorithm is explicit.
For a reduced fraction \(n/d\), \(d>1\), take
\(b=n^{-1}\bmod d\), \(0<b<d\), and
\(a=(nb-1)/d\). Then
\(\left(\begin{smallmatrix}-a&n\\-b&d\end{smallmatrix}\right)\)
has determinant one and gives the path from \(n/d\) to
\(a/b\). Continue from \(a/b\), whose denominator is
smaller, until an integer is reached. Polynomial
substitution gives the exact Manin coordinates. This is
implemented separately from the supplied continued-fraction
routine, and terminates by the strict denominator decrease.

In these period coordinates define
\[
 I_{n,m}(a)=\int_{a/5^n}^{i\infty}
             g_m(z)(5^nz-a)^w\,dz,
\qquad
 \mu_m(a+5^n\mathbb Z_5)
 =\alpha_m^{-n}I_{n,m}(a)
  -5^w\alpha_m^{-n-1}I_{n-1,m}(a).                    \tag{E.2.4}
\]
A fixed nonzero differential normalization is absorbed in
the period basis. Stabilizing the form as
\(g_m-(5^m/\alpha_m)g_m(5z)\) proves the distribution
relation: the substitution \(u=5z\) contributes the
factor \(1/5\) in the second term. All integral polynomial
paths have integral period coordinates; \(\alpha_m\) is
a unit. Hence (E.2.4) is a bounded integral two-coordinate
measure. The new program also verifies, over the integers,
\[
 \sum_{b=0}^4 I_{2,m}(a+5b)
 =a_5I_{1,m}(a)-5^w I_{0,m}(a),\qquad a=1,2,3,4.       \tag{E.2.5}
\]
For the twist it evaluates the exact polynomial paths
\[
 \sum_{v\bmod24}\chi_{24}(v)
 \int_{a/5^n+v/24}^{i\infty}
 g_m(z)(24\cdot5^nz-24a-5^nv)^w\,dz.                 \tag{E.2.6}
\]
The Gauss-sum identity identifies this with
\(24^w\tau(\chi_{24})I_{n,m,6}(a)\).
The factor is a five-adic unit, since
\(\tau(\chi_{24})^2=24\). Thus this exact twist formula
introduces no loss of precision or nonvanishing information.

For \(r=m+t\) put
\(J_{m,d,t}=\int_{\mathbb Z_5^\times}x^{-r}\,d\mu_{m,d}\).
On a unit residue class \(a+25\mathbb Z_5\),
\(x^{-r}\equiv a^{-r}\pmod{25}\). Bounded integrality
therefore proves that replacing the integrand by these
twenty values has error in \(25\mathbb Z_5^2\).
The values of \(\alpha_m\) in (E.2.2) decide each mass
modulo twenty-five exactly. The independently executed
outputs are
\[
\begin{array}{cc|cc}
m&t&d=1&d=6\\\hline
1&0&(22,3)&(19,6)\\
1&1&(0,5)&(0,5)\\
1&2&(19,6)&(24,1)\\
2&0&(9,9)&(9,9)\\
2&1&(6,19)&(6,19)\\
2&2&(7,7)&(7,7)\\
3&0&(20,15)&(4,18)\\
3&1&(16,0)&(7,0)\\
4&0&(5,0)&(5,0)\\
4&1&(0,20)&(0,20)\\
5&0&(6,3)&(12,6)\\
6&0&(12,12)&(12,12)
\end{array}                                                   \tag{E.2.7}
\]
All are nonzero. A vector divisible by five in this table
has valuation exactly one, so precision twenty-five is
essential for those entries. No unit claim is made for them.

The period sign has also been checked independently.
Reflection \(z\mapsto-\overline z\), including the
polynomial substitution \(P(X,Y)\mapsto P(-X,Y)\),
sends \(I_n(a)\) to \((-1)^wI_n(-a)\).
Its two-coordinate matrices are, in order \(m=1,\ldots,6\),
\[
 \begin{pmatrix}-1&2\\0&1\end{pmatrix},\quad
 \begin{pmatrix}0&-1\\-1&0\end{pmatrix},\quad
 \begin{pmatrix}1&0\\-6&-1\end{pmatrix},\quad
 \begin{pmatrix}-1&0\\0&1\end{pmatrix},\quad
 \frac1{173}\begin{pmatrix}-137&-30\\-372&137\end{pmatrix},\quad
 \begin{pmatrix}0&-1\\-1&0\end{pmatrix}.
\]
The program checks their squares, their action on every
Manin generator, and the sign \((-1)^{t+1}\) of every
moment in (E.2.7). Each eigenspace is one-dimensional.
Consequently each certificate concerns a single period
line, not a sum of unrelated periods. Coefficient duality
\((\operatorname{Sym}^wH)^*=
\operatorname{Sym}^wH\otimes\det(H)^{-w}\)
changes the homology sign by \((-1)^w\), giving the
cohomological period sign \((-1)^{m+t}\).

We now identify the specialization. For a primitive
character \(\eta\) of conductor \(5^n\), Fourier expansion
and the Gauss identity give
\[
 \sum_a\eta(a)^{-1}I_{n,m}(a)
 =\frac{(m-1)!\,5^{n(m-1)}\tau(\eta^{-1})}
 {(-2\pi i)^m}L(g_m,\eta,m).                           \tag{E.2.8}
\]
The old-level term of (E.2.4) sums to zero. Using
\(\tau(\eta)\tau(\eta^{-1})=\eta(-1)5^n\),
(E.2.8) matches the top-critical interpolation
\cite[Theorem 16.2, with \(r=m\)]{DE_Kato2004} at
\(\kappa^m\eta^{-1}\), up to a fixed nonzero period
scalar on each parity component. Rational period duality
allows these scalars as Kato's choices of periods.
Both transforms are bounded measures. A bounded power
series over a finite extension of \(\mathbb Q_5\) which
vanishes at all finite-order characters on a component
is zero, by the Weierstrass preparation theorem. Hence
the interpolation comparison identifies the whole
transform, and
\[
 J_{m,d,t}\ne0
 \quad\Longrightarrow\quad
 L_{p,\alpha_m}(g_{m,d})(\kappa^{-t})\ne0.              \tag{E.2.9}
\]
This proves the sign and the negative twist in the
specialization; neither is inferred just from the label
of a period vector.

Here is the exact part of Kato's theorem used. If \(T\)
is any stable lattice in the cohomological representation
of a good ordinary newform, and \(X(T)\) is the dual
ordinary cyclotomic Selmer module of \S17.3, then
\(X(T)\) is finitely generated torsion over
\(\Lambda=\mathbb Z_5[[\mathbb Z_5^\times]]\), and
\[
 \operatorname{length}_{\Lambda_{\mathfrak p}}X(T)_{\mathfrak p}
 \le \operatorname{ord}_{\mathfrak p}L_{p,\alpha}(g)
 \qquad(\operatorname{ht}\mathfrak p=1,\ 5\notin\mathfrak p).
                                                               \tag{E.2.10}
\]
This is \cite[Theorem 17.4(1)--(2)]{DE_Kato2004}.
Its assumptions contain no residual large-image condition.
The text at \S12.8 explicitly includes CM forms, with
the proof of Theorems 12.4--12.6 in \S15. The large-image
condition in 17.4(3) is not used. The ordinary filtration
and the identification of its discrete Selmer kernel are
17.1--17.2. All the hypotheses of (E.2.10) were checked
in (E.2.2). The official article metadata and the
reproduced primary theorem text are identified separately
in the source record; the inaccessible large PDF has not
been treated as if read.

Let \(\mathfrak p_t\) be the height-one evaluation prime
for \(\kappa^{-t}\), on its specified \(C_4\) component.
The nonzero value (E.2.9) has order zero at this prime;
this remains true for the entries of valuation one in
(E.2.7), because \(5\notin\mathfrak p_t\).
Thus (E.2.10) gives \(X(T)_{\mathfrak p_t}=0\).
Finite generation implies
\[
 \left(X(T)\otimes_{\Lambda,\kappa^{-t}}\mathbb Z_5\right)
 \otimes\mathbb Q_5=0.
\]
The integral tensor is finitely generated over
\(\mathbb Z_5\), hence finite. By contragredient
Pontryagin duality it is the dual of
\(\operatorname{Sel}_\infty(T)(-t)^{\Gamma}\).
Indeed the dual twist is \(t\), and taking coinvariants
imposes \(g=\kappa(g)^{-t}\) on \(X(T)\).
The relevant invariant subgroup therefore has bounded
exponent. This proves only the required specialization,
not vanishing of an entire Iwasawa branch.

We finish the descent to rational finite Selmer cohomology.
Let \(M=W_{m,d}(-t)\), \(T_M=T(-t)\), and \(B=M/T_M\).
The induced representation \(W_{m,d}\) is irreducible:
over \(K\) its two characters have distinct Hodge--Tate
weights, and complex conjugation exchanges them. If its
invariants under \(\Delta=G_{\mathbb Q(\mu_{5^\infty})}\)
were nonzero, normality and irreducibility would make the
whole representation invariant. It would then factor
through the abelian cyclotomic quotient. Complex
conjugation, with eigenvalues \(1,-1\), would commute
with the whole image and provide invariant lines, a
contradiction. Therefore \(M^\Delta=0\).

This proves finiteness of \(B^\Delta\) without a residual
irreducibility assumption. Choose finitely many elements
\(h_1,\ldots,h_a\in\Delta\) such that
\(M\to M^a\), \(v\mapsto((h_i-1)v)_i\), is injective;
such a finite choice exists by finite dimension and the
zero invariant subspace. Its inverse on its image has
bounded denominators relative to the lattices. Hence
\((h_i-1)v\in T_M\) for all \(i\) forces
\(v\in5^{-b}T_M\) for some fixed \(b\).
This bounds the order of every element of \(B^\Delta\)
and makes that group finite. Inflation--restriction now
bounds the exponent of the kernel of
\(H^1(\mathbb Q,B)\to H^1(\Delta,B)\).

All local finite conditions on \(M\) are zero. At five,
the de Rham filtration has \(\operatorname{Fil}^0\)
equal to the whole space, and there are no invariant
vectors, so \cite[Corollary 3.8.4]{DE_BlochKato1990}
gives \(H_f^1(\mathbb Q_5,M)=0\). The invariant assertion
follows either from the two CM characters or from the
nonunit complex absolute values of their Frobenius
roots. Away from five the representations have potential
good reduction and nonzero pure weight; Frobenius on
inertia invariants has no eigenvalue one. The unramified
finite local group is therefore zero. The real place
contributes zero at the odd prime five.

Use the fixed finite ramification set
\(\Sigma=\{2,3,5,\infty\}\). The group
\(H^1(G_{\mathbb Q,\Sigma},T_M)\) is finitely generated
and its image is a lattice in rational cohomology: finite
coefficient finiteness and compact Nakayama give the
first assertion, and every rational continuous cochain
on a compact group has bounded denominator, giving the
second. If \(H_f^1(\mathbb Q,M)\ne0\), its image under
reduction to \(B\) contains a nonzero divisible
five-primary group, namely that vector space modulo
its intersection with the integral lattice.
Every such reduced class is locally trivial. After
restriction to \(\Delta\) it satisfies the ordinary
kernel condition of Kato 17.2, and being restricted
from \(\mathbb Q\) it lies in
\(\operatorname{Sel}_\infty(T)(-t)^{\Gamma}\).
This target has bounded exponent by the specialization
argument. A divisible group has zero image in a group
of bounded exponent. It would therefore lie in the
bounded-exponent restriction kernel just proved, which
is impossible for a nonzero divisible group.
This contradiction proves (E.2.1) in every case.
\end{proof}

\begin{statement}{E.5}{Proved}{Specialization at a noncritical character without a pseudo-null gap}
Let \(k\) be imaginary quadratic and let \(p\) split in \(k\).
Let \(H\) be its Hilbert class field and let \(k_0/k\) be
abelian, contain \(H\), and satisfy
\(p\nmid[k_0:k]\#\mu(H)\). Let \(k_\infty/k\) be the
abelian extension with
\(\operatorname{Gal}(k_\infty/k_0)\simeq\mathbb Z_p^2\).
For any continuous one-dimensional character \(\rho\)
of \(\operatorname{Gal}(k_\infty/k)\), a nonzero Katz value
at \(\rho\) implies that the corresponding one-prime-ramified
Iwasawa module has zero rational \(\rho\)-specialization.
If \(\rho\ne1\) and \(V=\mathbb Q_p(\rho)\) has zero
finite local condition at every place, then
\(H_f^1(k,V)=0\).
\end{statement}
\begin{proof}
The original hypotheses are those of
\cite[\S4, pp.36--37]{DE_Rubin1991}. Theorem 4.1(i) gives
the split-prime main conjecture, in its elliptic-unit
form. The identification with the Katz characteristic
ideal is the Yager comparison, explicitly recorded in
\cite[Theorem 3.3.1]{DE_BCGKPST}. More than the
characteristic-ideal equality is needed for the desired
specialization. The additional precise theorem is
\cite[Theorem 5.3(v), pp.40--41]{DE_Rubin1991}:
the one-prime-ramified module \(X_{\mathfrak p}\) has no
nonzero pseudo-null submodule. In Rubin's notation
\(X(F)\) is defined on p.36 using extensions unramified
outside the chosen prime \(\mathfrak p\), not both
primes above \(p\). Theorem 5.3(iii) gives rank
\([k_{\mathfrak p}:\mathbb Q_p]-1=0\) in the split case.
These assertions apply to every finite-character
component, including the trivial one.

Here is the complete commutative-algebra step. Work in
one finite-character component of
\(\Lambda=\mathcal O[[\mathbb Z_p^2]]\), with coefficients
extended unramifiedly if needed. Let \(\mathfrak q\) be
the evaluation prime of \(\rho\). It has height two
and does not contain \(p\). The nonzero Katz value
makes the characteristic element a unit in
\(R=\Lambda_{\mathfrak q}\). Hence
\((X_{\mathfrak p})_{\mathfrak q}\) is zero at every
height-one prime of \(R\), and is a finite-length
\(R\)-module. If it were nonzero, the maximal ideal
of \(R\) would be associated to it. Localization of
associated primes would then make \(\mathfrak q\)
an associated prime of \(X_{\mathfrak p}\), giving
a submodule isomorphic to \(\Lambda/\mathfrak q\).
That is a nonzero pseudo-null submodule, contrary to
Rubin's Theorem 5.3(v). Thus
\((X_{\mathfrak p})_{\mathfrak q}=0\), which proves
zero rational specialization. Unramified faithfully
flat extension of coefficients preserves the
height and support assertions; alternatively the
argument descends to the characteristic element
over the finite coefficient ring before evaluation.
This is the missing justification for passing from
a nonzero characteristic value to a finite
specialization. It does not assume a pseudo-nullity
conjecture for the unramified Iwasawa module.

For the Selmer assertion choose a stable lattice \(T\subset V\)
and write \(B=V/T\). The rational global cohomology is finite
dimensional at a fixed finite ramification set, and the image
of integral cohomology is a full lattice. This follows from
finite-coefficient finiteness and compact Nakayama; every
continuous cocycle in \(V\) has bounded image and becomes
integral after multiplication by a power of \(p\).
A nonzero finite Selmer space would consequently give a
nonzero divisible subgroup of \(H^1(k,B)\) consisting of
locally trivial classes. Restriction to \(k_\infty\), where
\(B\) is constant, maps this subgroup into
\[
 \operatorname{Hom}_{\Delta\times\Gamma}
       (X_{\mathfrak p},B).
\]
The target is finite by the rational-specialization assertion
and finite generation. The restriction kernel has bounded
exponent: for an element \(g\) with \(\rho(g)\ne1\), Sah's
identity for the abelian group \(\Delta\times\Gamma\) kills
\(H^1(\Delta\times\Gamma,B)\) by \(\rho(g)-1\).
A nonzero divisible group cannot map into a finite group with
bounded-exponent kernel. This contradiction proves the
assertion. Faithfully flat extension of the coefficient ring
in the first part does not affect this vanishing.
\end{proof}

\begin{statement}{E.1}{Proved}{A second proof by twenty-four Eisenstein--Kronecker certificates}
For every \(d\in\{1,6\}\), \(1\le m\le6\), \(t\ge0\),
\(m+2t\le6\),
\[
 H_f^1\bigl(K,\mathbb Q_5(\psi_d^{-m-t}\overline\psi_d^{-t})\bigr)=0.
                                                               \tag{E.1.1}
\]
An independent computation from algebraic theta derivatives,
with a proved error in \(25\mathcal O\), gives the following
nonzero residues for the corresponding Katz specializations,
after multiplication by explicitly specified units:
\[
\begin{array}{cc|rr}
m&t&d=1&d=6\\\hline
1&0&7&24\\
1&1&20&15\\
1&2&6&9\\
2&0&11&11\\
2&1&14&14\\
2&2&12&12\\
3&0&5&9\\
3&1&11&16\\
4&0&10&10\\
4&1&3&3\\
5&0&11&22\\
6&0&23&23
\end{array}                                                    \tag{E.1.2}
\]
The proof uses Rubin's two-variable main conjecture and the
specialization argument E.5. It uses no modular-symbol
calculation and no theorem of Kato.
\end{statement}
\begin{proof}
We give the algebraic calculation, its precision, and the
arithmetic comparison separately. Write \(\varphi_d\) for
the algebraic Hecke character whose five-adic avatar is
\(\psi_d\), and put \(\mathfrak P=(1+i)\).
For a Gaussian integer \(\alpha\) prime to two, let
\(\varepsilon_1(\alpha)\in\{1,i,-1,-i\}\) be the unique unit
for which \(\varepsilon_1(\alpha)\alpha\equiv1\pmod{\mathfrak P^3}\).
Then
\[
 \varphi_1((\alpha))=\varepsilon_1(\alpha)\alpha,
 \qquad
 \varphi_6((\alpha))=\varepsilon_6(\alpha)\alpha,
 \quad
 \varepsilon_6(\alpha)=\varepsilon_1(\alpha)
                  \chi_{24}(\alpha\overline\alpha).             \tag{E.1.3}
\]
The first description follows equally from the CM Frobenius
formula or from the explicit theta series
\cite[\S3, pp.15--16]{DE_LiLongTu2018}. The second follows by
quadratic twisting by \(\sqrt6\).
Their conductors are \(\mathfrak P^3\) and \((12)\).
For the latter assertion (E.1.3) factors through
\((\mathbb Z[i]/12)^\times\). It is nontrivial at
\(7+6i\equiv1\pmod{\mathfrak P^3(3)}\), and at
\(1+4i\equiv1\pmod4\); both character values are \(-1\).
This excludes deletion of either the last two-adic conductor
step or the prime three. Enumeration gives
\(\#(\mathbb Z[i]/12)^\times=64\) and
\(\#\ker\varepsilon_6=16\). The program checks the character
law on all \(64^2\) pairs, and both conductor witnesses.
For even \(m\), the two target representations are equal,
since the quadratic twist occurs to the power \(m+2t\).

Use the invariant differential \(dx/(2y)\) on
\(y^2=x^3-d^2x\). Its period lattice is square. The reduced
theta function is its Weierstrass sigma function, because
\(e_2^*=0\). For a nonzero torsion point \(P\) prime to five,
put
\[
 X=\wp(P),\qquad Y=\wp'(P),\qquad
 Z=F_1(P)=\zeta(P)-\overline P/A_\Gamma.
\]
The algebra of its derivatives is
\[
 Y^2=4X^3-4d^2X,\quad
 DX=Y,\quad DY=6X^2-2d^2,\quad DZ=-X.                    \tag{E.1.4}
\]
For \(1\le k\le7\) set
\[
 c_1=Z,\qquad
 c_k=-\frac{D^{k-2}X}{k!}+\frac{G_k}{k}\ (k\ge2),\qquad
 G_4=d^2/15,
\]
with the other \(G_k\)'s in this range zero. Put \(e_0=1\)
and \(e_n=n^{-1}\sum_{k=1}^n k c_k e_{n-k}\).
Expansion of
\(\sigma(P+z+w)/(\sigma(P+z)\sigma(w))
\exp(-w\overline P/A_\Gamma)\) gives
\[
 \partial_z^A\partial_w^B\Theta_{P,0}(0,0)
       =D^A(B!e_{B+1})\qquad(A>0).                       \tag{E.1.5}
\]
Here the derivative means the regular coefficient after the
pole \(1/w\), whose \(z\)-derivative is zero. Equation
(E.1.5) is also the Eisenstein--Kronecker generating formula
\cite[Theorem 1.17]{DE_BannaiKobayashi2007}. Only \(k\le7\)
is needed below.

For \(d=1\) use
\[
 P=(i,1-i),\qquad 2P=(0,0),\qquad
 (X,Y,Z)=\bigl(i,2-2i,(1+i)/2\bigr).                    \tag{E.1.6}
\]
The displayed \(Z\) follows from the zeta duplication formula
and \(F_1(2P)=0\): \(Z=-(6X^2-2)/(4Y)\).
Moreover \([1+i]P=(1,0)\), so \(P\) generates
\(E_1[\mathfrak P^3]\). Its ray-class orbit has one element;
the sum over all residue units repeats it four times.
Changing the generator multiplies (E.1.5) by a fourth root
of unity, which has no effect on nonvanishing.

For the six odd-\(m\) entries with \(d=6\), work in the
algebra
\[
 R_0=(\mathbb Z/25)[r,v]/(r^2-2,\ v^8-6v^4-3).
                                                               \tag{E.1.7}
\]
Put \(s=(v^4-3)/2\), \(u=v^2/s\); then
\(s^2=3\), \(u^2=1+2s/3\). On \(E_6\) put
\[
 Q_3=(6u,12v),\qquad Q_4=(6(1+r),12s(1+r)),
 \qquad P=Q_3-Q_4.                                      \tag{E.1.8}
\]
These identities are initially interpreted in characteristic
zero and then reduced in (E.1.7). Direct substitution gives
\(3Q_3=0\), \(2Q_4=(6,0)\), \(3P=Q_4\), and \(4P=Q_3\).
For example the tangent at \(Q_4\) has slope \(s(2+r)\),
which gives its displayed double; the tangent at \(Q_3\)
has square \(18u\), so its double is \(-Q_3\).
As three is inert in \(K\), \(Q_3\) generates \(E_6[3]\).
The point \((6,0)\) is not killed by \(1+i\), so \(Q_4\)
has annihilator \((4)\). Thus \(P\) generates \(E_6[12]\)
as a \(\mathbb Z[i]\)-module.

CM reciprocity, in its explicit torsion form
\cite[Theorem 5.15 and Corollary 5.16]{DE_RubinAWS1999},
identifies the Galois orbit of this generator with
\(\ker\varepsilon_6\). Indeed multiplication by
\(\varphi_6((\alpha))=\varepsilon_6(\alpha)\alpha\)
runs through precisely that kernel, and ray-class
Chebotarev supplies each residue. Consequently
\([K(P):K]=16\). Conversely (E.1.8) recovers \(r,s,v\)
from \(P\), since \(3P=Q_4\), \(4P=Q_3\).
The rank-sixteen algebra in (E.1.7) is therefore the
reduction of the field algebra of this orbit. Its defining
polynomials have discriminants prime to five, so it is
\'etale at five. The trace is the orbit sum; on the displayed
basis it is exactly
\[
 \operatorname{Tr}\!\left(\sum_{a=0}^1\sum_{b=0}^7
                 z_{ab}r^av^b\right)=16z_{00}+48z_{04}. \tag{E.1.9}
\]
This follows from the Newton sums of \(v^8-6v^4-3\),
and from \(\operatorname{Tr}(r)=0\).
If \((X_j,Y_j)=(\wp(jP),\wp'(jP))\), the required
algebraic value of \(Z\) is
\[
 Z=\frac13\left[-\frac{6X_3^2-72}{4Y_3}
 -\frac{6X^2-72}{2Y}
 -\frac{Y_2-Y}{2(X_2-X)}\right].                        \tag{E.1.10}
\]
Apply duplication to \(3P\), whose double is a half-period,
and then the addition formula to \(2P+P\); this proves
(E.1.10). Every denominator used in (E.1.8)--(E.1.10)
is checked to be a unit in (E.1.7), with its inverse
verified by multiplication. Thus these formulas lose no
five-adic precision.

For completeness we identify the analytic moment and the
error bound. The integral theta model and a formal-group
isomorphism to \(\widehat{\mathbb G}_m\) give bounded
\(\mathcal O\)-valued measures \(\mu_{P,0}\), with a period
\(\Omega_5\in\mathcal O^\times\); this is
\cite[Corollary 2.17, Definition 3.2]{DE_BannaiKobayashi2007}.
The good integral model at five and the prime-to-five order
of \(P\) verify its hypotheses. The Yager ray-class sum
uses \(\mu_{P,0}(x,N(\mathfrak f)y)\). Its moments are
therefore the orbit sum of (E.1.5), multiplied by
\(N(\mathfrak f)^B\Omega_5^{A+B}\), and possibly by the
residue-unit multiplicity four. All these factors are units.
The restriction to both unit variables multiplies the orbit
sum by
\[
 \left(1-\frac{\pi^A}{\overline\pi^{B+1}}\right)
 \left(1-\frac{\pi^B}{\overline\pi^{A+1}}\right),
 \qquad\pi=-1+2i.                                     \tag{E.1.11}
\]
This follows directly by differentiating the four-term
restriction formula of
\cite[Proposition 3.5 and the proof of Theorem 3.7]{DE_BannaiKobayashi2007}.
Equivalently it is the two Euler factors in that theorem
with \(a=B\), \(b=A+1\). Both curves have
\(\varphi_d(\mathfrak p)=\pi\), since \(\chi_{24}(5)=1\).
Choose the embedding \(i\equiv18\pmod{25}\); then
\(\pi\equiv10\), \(\overline\pi\equiv13\).

The target moment is
\[
 \int_{(\mathbb Z_5^\times)^2}x^{-t-1}y^{m+t}
                                      \,d\mu_{\varphi_d}.
                                                               \tag{E.1.12}
\]
Take \(A=19-t\), \(B=m+t\). These are nonnegative,
\(A\ge17\), \(B\le6\), and
\(x^A\equiv x^{-t-1}\pmod{25}\) for every unit \(x\),
since the exponent of \((\mathbb Z/25)^\times\) is twenty.
Bounded integrality makes the error in replacing (E.1.12)
by this polynomial moment belong to \(25\mathcal O\).
The factors (E.1.11) are \(16\) modulo twenty-five when
\(B=1\), and one otherwise.
The exact rational differential algebra (E.1.4)--(E.1.5)
then gives (E.1.2). The first curve is evaluated in
\(\mathbb Q(i)\) before reducing, so cancellation of
factorial denominators is exact. For the second curve all
final derivative polynomials have integer coefficients;
this too is checked before reduction. The complete
executable computation is
\texttt{code/de\_ek\_certificates.py}, with its two derivative
modules and JSON outputs. It reads no modular-symbol data.
Every residue in (E.1.2) is nonzero, so its valuation is
exactly zero or one. An error in \(25\mathcal O\), or
multiplication by any of the specified units, cannot
change this nonvanishing conclusion.

We finally verify which character and which Iwasawa module
this computation concerns. For \(b>a\ge0\), the Yager
interpolation formula is proportional to
\[
 L_{\mathfrak f}(\overline\varphi_d^{a+b},b)
 =L_{\mathfrak f}((\varphi_d^b\overline\varphi_d^{-a})^{-1},0).
                                                               \tag{E.1.13}
\]
The proportionality factors are exactly its gamma, period
and Euler factors. Thus its measure is the Katz measure
on the character \(\rho=\psi_d^b\overline\psi_d^{-a}\),
with the shift \(x^{b-1}y^a\); see
\cite[Theorems 3.7 and 3.11]{DE_BannaiKobayashi2007}.
The identities on each finite-character branch determine
the bounded measures: these integer pairs are dense on
each such branch, and the integrands are continuous.
Consequently (E.1.12) is its specialization at
\(\rho=\psi_d^{-t}\overline\psi_d^{-m-t}\).
Omission of Euler factors at a larger fixed conductor
only multiplies the primitive Katz function by Euler
factors, so the nonzero imprimitive value implies a
nonzero primitive value. Conjugation in \(K/\mathbb Q\)
interchanges this representation with the one in (E.1.1)
and preserves its finite Selmer group.

Apply Rubin's two-variable theorem to the ray tower of
conductor \(\mathfrak f5^\infty\), with
\(\mathfrak f=\mathfrak P^3\) or \((12)\).
Here \(K\) has class number one, five splits, five is
prime to \(6\mathfrak f\), the finite ray extension at
\(\mathfrak f5\) has degree respectively sixteen or
256, and its pro-five Galois group has rank two.
It contains \(\mu_{5^\infty}\), by the norm character.
The character \(\rho\) factors through it by (E.1.3).
These check all splitting, conductor, prime-to-five and
finite-character hypotheses of Rubin's equality between
the Katz characteristic ideal and that of the one-prime
ramified module; the precise formulation is also
\cite[Theorem 3.3.1]{DE_BCGKPST}.
The equality \cite[Theorem 4.1(i)]{DE_Rubin1991} concerns each finite-character component,
including the trivial residual character. No big-image
or weight-less-than-five assumption is involved.

At the two primes above five the exponents of \(\rho\)
on the two inertia factors are \(-t\) and \(-m-t\).
At a prime where an exponent is zero, the remaining unramified
character has Frobenius a nonzero negative power of
\(\pi\) or \(\overline\pi\), and is not trivial.
The representation is crystalline with
\(D_{\rm dR}/\operatorname{Fil}^0=0\) and no local
invariants. Its finite local condition is zero by
\cite[Corollary 3.8.4]{DE_BlochKato1990}. Away from five
it is potentially unramified and pure of negative weight
\(-m-2t\), so Frobenius on inertia invariants has no
eigenvalue one; its unramified finite condition is zero.
There are no real places. Finally \(\rho\ne1\), since
\(m>0\). E.5 applies to every entry of (E.1.2) and gives
(E.1.1), with all restriction and lattice steps as proved
there. This completes the second proof.
\end{proof}

\begin{statement}{E.4}{Proved}{All vanishing defects and the dependent numerical tables}
None of the vanishings in Tasks E.1--E.3 fails.
Consequently all the defects
\(\delta_n(E_d)\), \(\delta_n^{\rm met}(E_d)\)
of LM8 F.9 are zero for \(3\le n\le8\).
The independent numerical recomputation gives,
for the full tower and \(s=0,1,2\), respectively,
\[
\begin{array}{cc|rrrrrrrr}
r&s&1&2&3&4&5&6&7&8\\\hline
0&0&0&0&1&3&6&11&20&36\\
0&1&0&1&2&4&7&12&21&37\\
0&2&0&2&3&5&8&13&22&38\\
1&0&1&1&2&4&7&12&21&37\\
1&1&1&2&3&5&8&13&22&38\\
1&2&1&3&4&6&9&14&23&39
\end{array}
\]
and for the metabelian tower,
\[
\begin{array}{cc|rrrrrrrr}
r&s&1&2&3&4&5&6&7&8\\\hline
0&0&0&0&1&3&5&8&11&15\\
0&1&0&1&2&4&6&9&12&16\\
0&2&0&2&3&5&7&10&13&17\\
1&0&1&1&2&4&6&9&12&16\\
1&1&1&2&3&5&7&10&13&17\\
1&2&1&3&4&6&8&11&14&18.
\end{array}
\]
The corresponding local dimensions are
\((1,2,4,7,13,22,40,70)\) and
\((1,2,4,7,11,16,22,29)\).
Thus the first strict levels remain one for \(r=0\)
and \(2,3,4\) for \(r=1\), \(s=0,1,2\).
This confirms the dependent entries for \(E_1\)
and \(E_6\) with the refined local conditions of
LM8 F.8--F.9.
\end{statement}
\begin{proof}
Every degree-\(n\) free-Lie constituent over \(K\)
has bidegree \((a,b)\) with \(a,b\ge1\), \(a+b=n\).
Its dual Tate twist is
\(\psi^{1-a}\overline\psi^{1-b}\).
For \(a>b\), put \(m=a-b\), \(t=b-1\).
Then \(m+2t=n-2\le6\), so E.1 applies; the
conjugate bidegree gives the same induced
representation. If \(a=b\), the rational
constituents are \(\mathbb Q_5(1-a)\) and its
quadratic CM twist, covered by E.3. The
metabelian layers are quotients of these
semisimple rational representations. All
stated defects therefore vanish.

We also verify the dimension calculation itself.
For the free Lie algebra on a two-dimensional
space with conjugation eigenvalues \(1,-1\),
Witt's character formula gives
\[
 r_n=\frac1n\sum_{d\mid n}\mu(d)2^{n/d},\qquad
 t_n=\frac1n\sum_{d\mid n}\mu(d)
          (\operatorname{Tr}(c^d))^{n/d},\qquad
 h_n^-=(r_n-t_n)/2,
\]
where \(\operatorname{Tr}(c^d)=0\) for odd \(d\)
and two for even \(d\).
For \(n=3,\ldots,8\), this gives
\(h_n^-=(1,2,3,5,9,16)\).
The metabelian layer is
\(\det(V)\otimes\operatorname{Sym}^{n-2}V\),
of dimension \(n-1\) and minus dimension
\(\lfloor n/2\rfloor\).
Poitou--Tate, with the now zero dual Selmer
spaces, gives exactly these global increments.
The higher central obstructions vanish as
well, since their rational \(H^2\) is dual
to those spaces. At non-five primes the
higher local cohomology vanishes by weights;
at five the whole higher \(H^1\) equals \(H_f^1\),
since \(\operatorname{Fil}^0L_n=0\) for
\(n\ge2\) and the dual invariants vanish
for \(n\ge3\). Thus the increments impose
no additional higher local equation.
The initial refined components have dimensions
\(r\) and \(r+s\): the class-two central
\(\mathbb Q_5(1)\) coordinate is the unit
valuation coordinate, with one free
parameter for each allowed denominator
prime. Adding the increments proves the
two tables. Locally degree one contributes
one, and every higher degree contributes
its full rank, giving the two local rows.
The executable integer calculation
\texttt{code/de\_tables.py} independently
computes all ranks, traces, sums and first
strict levels from these formulas.
\end{proof}
